\metatitle{Littlewood--Richardson coefficients}

\metadescription{Various models for Littlewood--Richardson coefficients}
\metakeywords{Littlewood--Richardson coefficient, Littlewood--Richardson
tableau, hive, Knutson--Tao puzzle, Berenstein--Zelevinsky pattern,
stretched Littlewood--Richardson coefficient}


\section[littlewoodRichardsonModels]{Combinatorial models for Littlewood--Richardson coefficients}

The \hyperref[schurLittlewoodRichardson]{Littlewood--Richardson coefficients}
are defined as the structure constants $c^{\lambda}_{\mu\nu}$ appearing in the
expansion
$\schurS_\mu \schurS_\nu =
\sum_\lambda c^{\lambda}_{\mu\nu} \schurS_\lambda$.
One can show that $c^{\lambda}_{\mu\nu}=0$ unless $\mu,\nu \subseteq \lambda$.
Furthermore, we have the symmetries
\[
c^{\lambda}_{\mu\nu} = c^{\lambda}_{\nu\mu} = c^{\lambda'}_{\mu'\nu'}.
\]
A simple recursion for computing $c^{\lambda}_{\nu\mu}$
can be obtained from the \hyperref[schurShiftedLittlewoodRichardson]{shifted Schur functions}.

The Littlewood--Richardson coefficients \hyperref[kostkaFromLR]{generalize the skew Kostka coefficients}.
\name{Igor Pak} and \name{Ernesto Vallejo} compare several standard models
for these coefficients as integer points in Littlewood--Richardson cones
\cite{PakVallejo2005}.  They give explicit linear bijections between
Littlewood--Richardson tableaux, Knutson--Tao hives, and
Berenstein--Zelevinsky triangles, showing that these models have the same
underlying combinatorial structure.
The \hyperref[hivePolytopes]{hives and BZ-polytopes page} collects the
polyhedral viewpoint, saturation theorem, and Ehrhart stretching questions.

\hyperref[jeuDeTaquinOperation]{Jeu-de-taquin} gives a formulation using
standard tableaux.  If $U$ is a fixed standard Young tableau of shape $\nu$,
then $c^\lambda_{\mu\nu}$ is the number of standard tableaux of skew shape
$\lambda/\mu$ whose rectification is $U$ \cite[Sec.~5]{Fulton1997}.  This
count is independent of the chosen standard tableau $U$ of shape $\nu$.


\name{Olga Azenhas} proves that for fixed $\lambda$ and $\mu$, the set
\[
 \{ \nu : c^{\lambda}_{\mu \nu}>0 \}
\]
is a sub-lattice of the partition lattice under domination order
\cite{Azenhas1999}.
See also \cite{TaskinGumus2022x} for conditions on when
$c^{\lambda}_{\mu \nu}$ is positive.


The \defin{Horn inequalities}\label{hornInequalities} determine for which
triples $(\lambda, \mu, \nu)$ we have $c^{\lambda}_{\nu\mu} > 0$; see
\cite{Klyachko1998,KnutsonTau1999}.



\subsection[lrTableaux]{Littlewood--Richardson tableaux}

The Littlewood--Richardson rule states that $c^{\lambda}_{\mu\nu}$ counts the
number of skew \hyperref[prelimTableaux]{semistandard Young tableaux} of shape
$\lambda/\mu$ with content $\nu$, with the additional restriction that the
concatenation of the reversed rows is a
\defin{lattice word}\label{lrLatticeWord}.
A word $w_1,\dotsc,w_\ell$ is a lattice word if every prefix has the property
that its content is a partition.  That is, the number of times $i$ appears in
the prefix is at least as many as $i+1$, for every $i$.
Tableaux of this type are called
\defin{Littlewood--Richardson tableaux}\label{littlewoodRichardsonTableaux}.

\begin{example*}[Littlewood--Richardson tableaux]
Consider the coefficient of $\schurS_{542}$ in the product $\schurS_{32}\cdot \schurS_{321}$.
We must find all Littlewood--Richardson tableaux of shape $542/32$ with content $321$.
There are two such tableaux:
\begin{figure}
\begin{ytableau}
\none & \none & \none & 1 & 1 \\
\none & \none & 1 & 2 \\
2 & 3
\end{ytableau}
\begin{ytableau}
\none & \none & \none & 1 & 1 \\
\none & \none & 2 & 2 \\
1 & 3
\end{ytableau}
\end{figure}
Consequently, $c^{542}_{32,321}=2$.
\end{example*}


\subsection[lrYamanouchi]{Semistandard Young tableaux}

Lothaire gives the following characterization of Littlewood--Richardson
coefficients \cite[Cor.~5.4.8]{Lothaire2002}:
$c^{\lambda}_{\mu\nu}$ is the number of
\hyperref[prelimTableaux]{semistandard Young tableaux} $T$ with shape $\nu$
and content $\lambda-\mu$ such that
$\rw(t) \cdot \ell^{\mu_\ell} \dotsm 2^{\mu_2}1^{\mu_1}$ is a Yamanouchi word.
Here \hyperref[yamanouchiWord]{Yamanouchi word} means that every prefix has
partition content.

Using \cite[Rem.~14 and Cor.~21]{KaliszewskiMorse2019}, one can show that
this is equivalent to the following model.


\subsection[lrMolevSagan]{Semistandard Young tableaux II}

\name{Sergey Fomin}, \name{Curtis Greene}, \name{Alexander Molev}, and
\name{Bruce Sagan} describe the following model
\cite{FominGreene1993,MolevSagan1999}.
$c^{\lambda}_{\mu\nu}$ is the number of
\hyperref[prelimTableaux]{semistandard Young tableaux} $T$ with shape $\mu$
and with the property that the column index reading word of $T$
\defin{fits}\label{lrFitsSkewShape} the skew shape $\lambda/\mu$.
That is, we take the column reading word of $T$, reading columns right to left.
This gives a word $w$. For each entry $w_i$, we put an entry $i$ in row $w_i$,
so that the resulting filling $F=F(T)$ is of shape $\lambda/\mu$.
If this filling is a skew SYT, then we say that $T$ fits $\lambda/\mu$.

Note that the conditions force that $T$ has content $\lambda-\mu$.

Observe that this model is generalized to Littlewood--Richardson coefficients for
\hyperref[schurShifted]{shifted Schur functions}.

\name{Chau Nguyen}, \name{Son Nguyen}, and \name{Dora Woodruff} give a
formula for Littlewood--Richardson coefficients using peelable tableaux
compatible with shuffle tableaux \cite{NguyenNguyenWoodruff2025x}.  Their
rule behaves well with Bender--Knuth involutions, recovers the usual symmetry
of Littlewood--Richardson coefficients, and proves a special case of the
Lam--Postnikov--Pylyavskyy Schur log-concavity conjecture.


\begin{example*}[Fomin--Greene model]
Let us show that $c^{542}_{321,32}=2$ using the Fomin--Greene model.
We are counting SSYT with shape $\mu=321$,
which fit $542/32$.

There are two such SSYT, together with the corresponding SYT obtained from
the column reading word:
\begin{figure}
\begin{ytableau}
1 & 1 & 2\\
2 & 3 \\
3
\end{ytableau}
\begin{ytableau}
\none & \none & \none & 2 & 4 \\
\none & \none & 1 & 5 \\
3 & 6
\end{ytableau}
\end{figure}

\begin{figure}
\begin{ytableau}
1 & 1 & 3\\
2 & 2 \\
3
\end{ytableau}
\begin{ytableau}
\none & \none & \none & 2 & 4 \\
\none & \none & 3 & 5 \\
1 & 6
\end{ytableau}
\end{figure}
Consequently, $c^{542}_{32,321}=2$.
\end{example*}


\subsection[lrCharge]{Postfix-charge model}

In joint work with \name{Joakim Uhlin}, we show
\cite{AlexanderssonUhlin2020} that
\[
c^{\lambda}_{\mu\nu}
=
\left|
 T \in \SSYT(\nu, \lambda - \mu) :
\charge( \rw(T) \cdot \ell^{\mu_\ell} \dotsm 2^{\mu_2}1^{\mu_1})=0
\right|.
\]
That is, $c^{\lambda}_{\mu\nu}$ is equal
to the number of \hyperref[prelimTableaux]{semistandard Young tableaux}
of shape $\nu$ and content $\lambda - \mu$ such that the reading word
postfixed with $\ell^{\mu_\ell} \dotsm 2^{\mu_2}1^{\mu_1}$ has
\hyperref[chargeStatistic]{charge} equal to zero.


\begin{example*}[Postfix-charge model]
Let us compute $c^{33221}_{221,321}$, where
$\lambda = 33221$, $\mu = 221$ and $\nu = 321$.
We are interested in semistandard tableaux with shape $321$ and content
$\lambda-\mu = 11121$.
There are eight such tableaux.
\begin{figure}
\ytableaushort{144,25,3}\ytableaushort{135,24,4}\ytableaushort{134,25,4}\ytableaushort{134,24,5}
\ytableaushort{125,34,4}\ytableaushort{124,35,4}\ytableaushort{124,34,5}\ytableaushort{123,44,5}
\end{figure}
We take their reading words and concatenate
$\ell^{\mu_\ell} \dotsm 2^{\mu_2}1^{\mu_1} = 32211$ at the end.
\begin{array}{ll}
\toprule
\text{Word} & \text{Charge} \\
\midrule
 32514432211 & 1 \\
 42413532211 & 1 \\
 42513432211 & 0 \\
 52413432211 & 0 \\
 43412532211 & 2 \\
 43512432211 & 1 \\
 53412432211 & 2 \\
 54412332211 & 1 \\
\bottomrule
\end{array}
As we see, the third and fourth tableaux above give charge $0$, so
$c^{33221}_{221,321}=2$.
\end{example*}

\begin{example*}[Postfix-charge model II]
Let us compute $c^{542}_{32,321} = c^{\lambda}_{\mu\nu}$.
We are interested in semistandard tableaux with shape $321$ and content
$\lambda-\mu = 222$.
There are two such tableaux.
\begin{figure}
\ytableaushort{113,22,3}
\ytableaushort{112,23,3}
\end{figure}
We take their reading words and concatenate $22111$ at the end,
$322113\;22111$ and $323112\;22111$.
Both of these words have charge $0$, so $c^{542}_{32,321}=2$.
\end{example*}

%
% Table[
%  inSchur =
%   ToSchurBasis[
%    SchurSymmetric[ConjugatePartition@nu] SchurSymmetric[
%      ConjugatePartition@mu], ss];
%  Table[
%   lrCoeff = Coefficient[inSchur, ss[ConjugatePartition@lam]];
%   wt = PadRight[lam, 7] - PadRight[mu, 7];
%   ssyts = SemiStandardYoungTableaux[{nu, {}}, wt];
%   postFix =
%    Flatten@Reverse@MapIndexed[ConstantArray[#2[[1]], #1] &, mu];
%   lrCoeff2 =
%    Length@Select[ssyts,
%      WordCharge[Join[SYTReadingWord@#, postFix]] == 0 &];
%   {lam, mu, nu} -> {lrCoeff, lrCoeff2};
%   lrCoeff2 == lrCoeff
%   , {lam, IntegerPartitions[7]}]
%  , {nu, IntegerPartitions[4]},
%  {mu, IntegerPartitions[3]}]




\subsection[bzPatterns]{Berenstein--Zelevinsky patterns}

\name{Arkady Berenstein} and \name{Andrei Zelevinsky}
\cite{BerensteinZelevinsky1992} gave a combinatorial model for the
Littlewood--Richardson coefficients that uses certain triangular arrays,
reminiscent of \hyperref[gtpatterns]{Gelfand--Tsetlin patterns}.
The exact definition is written in an accessible manner in
\cite[Def.~A1.3.10]{StanleyEC2}.


\subsection[knutsonTaoHoneycomb]{Knutson--Tao honeycombs}


\subsection[knutsonTaoWoodwardPuzzles]{Knutson--Tao--Woodward puzzles}

Knutson--Tao--Woodward puzzles
\cite{KnutsonTau1999,KnutsonTaoWoodward2004} are another model used for
proving the saturation conjecture.
An overview of the saturation conjecture is given by \name{Anders Buch}
\cite{Buch2000}.

\name{Christian Gaetz} and \name{Yibo Gao} relate Schubert puzzles to
interlacing triangular arrays \cite{GaetzGao2025}.  They prove that certain
rank-three interlacing arrays are equinumerous with proper colorings of a
triangular grid graph via generalized Knutson--Tao puzzles, and use the same
objects to compute structure constants in $K(\operatorname{Gr}(d,n))$, in the
localized cohomology of cotangent bundles of \hyperref[grassmannian]{Grassmannians},
and in the cohomology of partial \hyperref[flagVarieties]{flag varieties}.

\name{Portia Anderson} extends triangular Schubert puzzles to puzzles with
convex polygonal boundary shapes \cite{Anderson2024Puzzlesx}.  The resulting
four-, five-, and six-sided puzzles have
\hyperref[schubertCalculus]{Schubert-calculus} interpretations,
give commutativity statements for several symmetric boundary conditions, and
express the associated structure constants in terms of
Littlewood--Richardson numbers.

\name{Emmanuel Briand} and \name{Mercedes Rosas} study the full group of
symmetries of the Littlewood--Richardson coefficients for $SL_3$
\cite{BriandRosas2020x}.  \name{Maxime Pelletier} and
\name{Nicolas Ressayre} exhibit an unexpected hidden symmetry of
Littlewood--Richardson coefficients \cite{PelletierRessayre2022}; a proof is
given by \name{Darij Grinberg} \cite{Grinberg2021}.
\name{Olga Azenhas}, \name{Alessandro Conflitti}, and \name{Ricardo Mamede}
describe a faithful action of $\setZ_2\times D_3$ on Littlewood--Richardson
tableaux, companion tableaux, hives, and puzzles \cite{AzenhasConflittiMamede2025x}.
This gives a uniform way to relate commutativity, conjugation, and cyclic
symmetries of Littlewood--Richardson coefficients.
\name{David E. Speyer} proves the conjecture of
\name{Thomas Lam}, \name{Alexander Postnikov}, and \name{Pavlo Pylyavskyy}
on Schur nonnegativity of certain differences
$\schurS_{\lambda'}\schurS_{\mu'}-\schurS_\lambda\schurS_\mu$
\cite{Speyer2026x}.  The proof introduces a new model for
Littlewood--Richardson coefficients, called skeps, and proves an
L-log-concavity theorem using Murota's L-convexity.




\subsection[steinbergLRFormula]{Steinberg's formula}

One can express the
\hyperref[schurLittlewoodRichardson]{Littlewood--Richardson coefficients}
as a signed sum involving the
\hyperref[kostantPartitionFunction]{Kostant partition function}, by using
Steinberg's formula \cite{Steinberg1961}.

See the exact formula in the section with the
\hyperref[littlewoodRichardsonFromKostant]
{Littlewood--Richardson coefficients from the Kostant partition function}.



\subsection[socleTableaux]{Socle tableaux}


\name{Justyna Kosakowska} and \name{Markus Schmidmeier} realize
$c^{\lambda}_{\mu\nu}$ as the number of \emph{socle tableaux} of shape
$\lambda/\nu$ and content $\mu$ \cite{KosakowskaSchmidmeier2022}.
Note that the roles of $\mu$ and $\nu$ have been interchanged.

A \defin{socle tableau}\label{socleTableau} is weakly decreasing in rows,
strictly decreasing in columns, and for each vertical line and integer $\ell$,
at most as many entries $\ell+1$ occur on the left-hand side as entries
$\ell$.


\subsection[lrExplicitFormula]{Evaluating Schur functions at roots of unity}

\name{Xueqing Wen} gives a fairly explicit formula
for Littlewood--Richardson coefficients, as a sum over partitions,
where the terms include a product of three Schur functions evaluated at
certain roots of unity; see \cite{Wen2021x}.



\section[newellLittlewood]{Newell--Littlewood numbers}

The \defin{Newell--Littlewood numbers}\label{newellLittlewoodNumbersDefinition}
may be defined as
\[
 N_{\mu,\nu,\lambda}
 =
 \sum_{\alpha,\beta,\gamma}
 c^{\mu}_{\alpha\beta} c^{\nu}_{\alpha\gamma}
 c^{\lambda}_{\beta\gamma}.
\]
For recent papers on these coefficients, see \cite{GaoOrelowitzYong2020x,GaoOrelowitzYong2020IIx}.

The Littlewood--Richardson coefficients arise as tensor product
multiplicities of irreducible \hyperref[representationTheory]{representations}
of $\GL(V)$.
The Newell--Littlewood numbers arise in the same manner, where $\GL(V)$ is
replaced by a Lie group of type $B$, $C$, or $D$.

There is a symmetric function basis, the
\defin{Koike--Terada basis}\label{koikeTeradaBasis},
$\{\schurS_{[\lambda]}\}_\lambda$,
for which the $N_{\mu,\nu,\lambda}$ serve as multiplication structure constants.
\name{Shiliang Gao}, \name{Gidon Orelowitz}, and \name{Alexander Yong} give
an explicit Jacobi--Trudi type formula for this basis
\cite{GaoOrelowitzYong2020x}.

They also give a polytope description for $N_{\mu,\nu,\lambda}$ and state a
saturation-type conjecture \cite[Thm.~5.1]{GaoOrelowitzYong2020x}.  See also
\cite{King2021x} for more conjectures.  The saturation conjecture for the
Newell--Littlewood numbers has since been proved in \cite{Min2024x}.

\section[lrStretched]{Stretched Littlewood--Richardson coefficients}

The Littlewood--Richardson coefficients are in correspondence with lattice
points in certain rational \hyperref[polytope]{polytopes}, as proved in
\cite{BerensteinZelevinsky1988}.
Consequently, it is immediate that this function is a quasipolynomial, since
we can model it as the \hyperref[ehrhartPolynomialDefinition]{Ehrhart}
quasipolynomial of a rational polytope.
Polynomiality was later proved by \name{Harm Derksen} and
\name{Jerzy Weyman} \cite{DerksenWeyman2002}, and by
\name{Etienne Rassart} \cite{Rassart2004}.
\name{Warut Thawinrak} gives a simple proof of polynomiality
\cite{Thawinrak2022x}.
The proof uses the \hyperref[kostantPartitionFunction]{Kostant partition function}
formula due to \name{Robert Steinberg}.




\begin{conjecture}[See \cite{KingTolluToumazet2004}]
The function
\[
t \to c^{t \nu}_{t\lambda,t\mu}
\]
for fixed partitions is a polynomial in $t$.
Note that this is the \hyperref[ehrhartPolynomialDefinition]{Ehrhart polynomial}
of the corresponding \hyperref[bzPatterns]{BZ-polytope}, where lattice
points are BZ-patterns.

Computer experiments in \cite{KingTolluToumazet2004} suggest that the Ehrhart
polynomial above has nonnegative coefficients.
This conjecture is still open.
A special case of this conjecture is to consider the map $t \to K_{t\lambda,t\mu}$,
of so called \hyperref[kostkaStretched]{stretched Kostka coefficients}.
Various aspects of this problem have been considered in
\cite{Alexandersson2016GTPoly,Alexandersson2019CounterExamples}.
\end{conjecture}

\defin{Fulton's conjecture}\label{fultonsConjecture} stated that
$c^{\nu}_{\lambda,\mu}=1$ implies that $c^{t \nu}_{t\lambda,t\mu}=1$ for all
$t \geq 1$.  This was later proved in \cite{KnutsonTaoWoodward2004}.

\name[Mrigendra Singh Kushwaha]{Mrigendra Singh Kushwaha},
\name[K. N. Raghavan]{K.~N. Raghavan}, and
\name[Sankaran Viswanath]{Sankaran Viswanath}
study refined Littlewood--Richardson coefficients attached to
\hyperref[key]{key polynomials} \cite{KushwahaRaghavanViswanath2021}.
If $\chi_{w,\mu}$ is a Demazure character, then the coefficient
$c^\nu_{\lambda\mu}(w)$ is defined by
\[
T_{w_0}\left(x^\lambda \chi_{w,\mu}(\xvec)\right)
= \sum_\nu c^\nu_{\lambda\mu}(w)\schurS_\nu(\xvec).
\]
They give a hive model for these coefficients and prove a saturation theorem
when $w$ is $312$-avoiding, $231$-avoiding, or a commuting product of such
permutations, generalizing the Knutson--Tao saturation theorem.

The next natural case is $c^{\nu}_{\lambda,\mu}=2$.
It was proved by \name{Christian Ikenmeyer} \cite{Ikenmeyer2016} and later by
\name{Cass Sherman} \cite{Sherman2015x,Sherman2017} that
\[
c^{\nu}_{\lambda,\mu}=2 \implies c^{t \nu}_{t\lambda,t\mu}=t+1.
\]

\name{Pierre-Emmanuel Chaput} and \name{Nicolas Ressayre} give a nice
generalization of this implication \cite{ChaputRessayre2025}:
\[
c^{\nu}_{\lambda,\mu}=2 \implies c^{\nu(s,t)}_{\lambda(s,t),\mu(s,t)}=\binom{s+t}{t},
\]
with the convention that $\lambda(s,t) \coloneqq s (t \lambda')'$; that is,
we stretch in both the vertical and horizontal directions.


Stretching of Newell--Littlewood coefficients is studied in \cite{King2021x}.
This preprint contains many conjectures analogous to results and conjectures
about the Littlewood--Richardson coefficients.  \name{Ronald C. King} proves
that the map $t \to n^{t \nu}_{t\lambda,t\mu}$ is a quasipolynomial with
period at most $2$.


\section[lrIdentities]{Identities}


\subsection[lrCoquereauxZuber]{Coquereaux--Zuber relation}

Let $\mu, \nu \vdash n$ be such that all entries are at most $k$.
\name{Robert Coquereaux} and \name{Jean-Bernard Zuber} prove that
\[
 \sum_{\lambda \vdash n} c^{\lambda}_{\mu, \nu} =
 \sum_{\lambda \vdash n} c^{\lambda}_{\mu^{\vee k}, \nu},
\]
where $\mu^{\vee k} \coloneqq (k-\lambda_n,\dotsc,k-\lambda_2,k-\lambda_1)$
\cite{CoquereauxZuber2011}.



\subsection[harrisWillenbring]{The Harris--Willenbring identity}

See \cite{HarrisWillenbring2014}.

\[
 \sum_{n \geq 0}
 \sum_{\lambda,\mu, \nu \vdash n}
 \left( c^{\lambda}_{\mu\nu} \right)^2
 t^{|\mu|}q^{|\nu|}
 =
 \prod_{i \geq 1} \frac{1}{1-t^i-q^i}
\]



\subsection[lrShrivastava]{Shrivastava's relation with Kostka numbers}

\name{Sagar Shrivastava} shows how one can compute
Littlewood--Richardson coefficients as a signed sum of
\hyperref[kostkaCoefficient]{Kostka coefficients} \cite{Shrivastava2022x}.

It is known that (skew) \hyperref[kostkaCoefficient]{Kostka coefficients}
can be obtained as a
\hyperref[kostkaFromLR]{special case of Littlewood--Richardson coefficients};
see, for example, \cite[p.~338, Eq.~7.61]{StanleyEC2}.




\section[lrcalc]{Computing coefficients, and A. Buch's lrcalc}

\name{Hariharan Narayanan} proved in 2006 that computing
Littlewood--Richardson coefficients and
\hyperref[kostkaCoefficient]{Kostka coefficients} is in general
$\#P$-complete \cite{Narayanan2006}.
This rules out efficient algorithms for computing these quantities
(assuming widely believed conjectures regarding complexity theory).

\name[Peter Bürgisser]{Peter Bürgisser} and \name{Christian Ikenmeyer}
give a polynomial-time algorithm for deciding whether
$c^\nu_{\lambda,\mu}>0$ \cite{BurgisserIkenmeyer2013}.  Starting from the
Knutson--Tao hive model, they interpret $c^\nu_{\lambda,\mu}$ as the number
of capacity-achieving hive flows on the honeycomb graph and obtain an
$O(n^3\log \nu_1)$ positivity test.

\name{Anders Buch} has a fast C program, \texttt{lrcalc}, for computing
Littlewood--Richardson coefficients \cite{BuchLRCalc}.

See also \name{Matthew J. Samuel}'s \texttt{schubmult} Python package for
computing products of \hyperref[schubert]{Schubert polynomials}, which in
particular can compute Littlewood--Richardson coefficients
\cite{SamuelSchubmult}.
